\chapter{Fold-level investigation: what actually goes wrong}
\label{ch:folds}

\begin{saybox}[title={What to say at the start of this section}]
This is the section that answers ``so what is the model actually bad at?''
Do not speak in generalities about distribution shift. Name the rule, name
the cheese, name the preservative, give the numbers in days. Every claim in
this chapter is traceable to a specific set of rows in
\texttt{rf\_fold\_predictions.csv}.
\end{saybox}

\section{Method for this chapter}

The historical leave-rule-out experiments stored only per-fold summary
statistics. To investigate individual folds, the evaluation was re-executed
with identical data, identical hyperparameters and the identical fold
construction, retaining every out-of-fold prediction
(\evd{scripts/experiments/leave\_rule\_out\_row\_level.py}).

\begin{limitbox}[title={Reconstruction, clearly labelled}]
The row-level predictions analysed here come from a \textbf{re-execution},
not from the original diagnostic run. The original summary files are
retained unmodified in
\evd{scripts/experiments/overfitting\_diagnosis/}. The reconstruction
reproduces the original per-fold $\rsq$ values to the precision reported,
which is expected given the shared seed and the deterministic fold
construction, but the distinction is preserved rather than glossed: the
figures in this chapter are labelled as reconstructed.
\end{limitbox}

Each of the $56$ folds was predicted by all three tree ensembles, giving
$168$ fold results and $102\,000$ row-level predictions with full
formulation context attached: cheese product, physical form, treatment
ingredient and concentration, storage temperature, packaging, indicator
type and threshold, the observed target, the prediction, and the signed
error.

\tbllong{tab_folds_random_forest}{Complete Random Forest fold results. Every
held-out generation rule in every specialist, sorted by $\rsq$ within
specialist. Negative $\rsq$ values are bold.}

\section{The four weakest folds}

\tbl{tab_worst_fold_detail}{Side-by-side diagnostic profile of the four
weakest Random Forest folds. Every cell is computed from the held-out rows
themselves or from the corresponding training partition.}

\fig{fig31_worst_folds_scatter}{0.98}{%
  Actual against predicted shelf life for the four weakest folds. Note the
  scale differences: the two natamycin folds on the right are compressed
  into a few days, while the clostridia fold on the left spans hundreds.}

\fig{fig21_fold_size_vs_r2}{0.88}{%
  Fold $\rsq$ against the number of rows in the held-out rule for all
  $168$ fold results. Large folds are consistently well predicted; the
  dispersion is almost entirely among rules with fewer than $200$ rows ---
  with one crucial exception at $620$ rows, analysed below.}

\section{Fold 1: \texttt{CLOSTRIDIA|CLOSTRIDIA\_GRANA|HARD\_SEMI\_REVIEW}}
\label{sec:fold-clostridia}

This is the worst fold in the entire study and the only large one that
fails. It is also the most instructive, because its cause can be
established rather than hypothesised.

\subsection{What was held out}

\begin{description}[leftmargin=1.5em, style=nextline, font=\bfseries\color{dgreen}, itemsep=1pt]
\item[Specialist] hard / general shelf life
\item[Rows] $620$ held out, $3\,580$ remaining for training
\item[Cheese products] pecorino romano, gruyère, manchego, parmigiano
  reggiano, cheddar, emmental --- \emph{all six also present in training}
\item[Treatments] none (control), lysozyme, potassium nitrate, protective
  \emph{Lactobacillus} culture
\item[Indicator] \texttt{clostridia} --- exclusively
\item[Storage temperature] $8$--$15^\circ$C (training spans $2$--$15^\circ$C)
\item[Ripening] $78.5$--$843.1$ days (training spans $67.2$--$877.9$)
\end{description}

\subsection{Observed result}

Random Forest: $\rsq = -3.860$, MAE $= 67.4$ days, RMSE $= 83.1$ days, mean
signed error $= -62.2$ days. The model under-predicts by roughly two months
on average. Actual targets: mean $189.0$ days, range $69.5$--$240.0$,
$s_y = 37.7$. Predictions: mean $126.8$ days.

\subsection{Diagnosis}

Working through the candidate explanations in order:

\paragraph{Is it an unseen product?} No. All six cheese products in the
held-out rule also appear in the training partition. Product novelty is
eliminated.

\paragraph{Is it fold size?} No. At $620$ rows this is the third-largest
fold in the specialist. Small-sample instability is eliminated.

\paragraph{Is it low target variance?} No. $s_y = 37.7$ days is
substantial. The $\rsq$ denominator is large; the numerator is simply
larger.

\paragraph{Is it an unseen \emph{mechanism}?} Yes --- and this is
demonstrable. Cross-tabulating \texttt{indicator\_type} against
\texttt{source\_rule\_id} within hard/general shows that the
\texttt{clostridia} indicator appears in exactly two rules: this one
($620$ rows) and a small sibling,
\texttt{CLOSTRIDIA|CLOSTRIDIA\_GRANA|HARD\_SEMI\_REVIEW|SHRED\_NAT}
($32$ rows). When this rule is withheld, only $32$ clostridia rows remain
in training out of $3\,580$.

The clostridia endpoint operates on a completely different time scale from
every other indicator in the specialist:

\begin{center}\small
\begin{tabular}{lrrr}
\toprule
\texttt{indicator\_type} & rows & mean $y$ (d) & range (d) \\
\midrule
\textbf{clostridia} & 652 & \textbf{188.8} & 69.5--240.0 \\
molds & 432 & 68.3 & 23.6--132.5 \\
oxidation\_index & 876 & 83.2 & 27.1--162.3 \\
peroxide\_value & 388 & 81.3 & 35.5--134.8 \\
ph & 356 & 81.5 & 32.7--155.9 \\
sensory\_acceptability & 1\,076 & 81.9 & 31.1--177.7 \\
total\_viable\_count & 988 & 75.1 & 19.6--147.3 \\
yeasts\_molds & 1\,232 & 67.2 & 17.3--137.7 \\
\bottomrule
\end{tabular}
\end{center}

Late blowing by clostridia is a slow ripening-associated defect; the other
indicators track surface or bulk spoilage that manifests in weeks. The
mean target for clostridia rows is $2.3\times$ that of any other indicator.

\paragraph{The mechanism of failure.} With the rule withheld, the training
target mean falls to $76.9$ days against the held-out mean of $189.0$. By
\cref{eq:rf-bound}, the forest's prediction is a weighted average of
training targets, so it is structurally confined near the training range.
It predicts $126.8$ days on rows whose truth averages $189.0$ --- a
$-62.2$ day bias that is not a tuning failure but a consequence of what
tree averaging can represent.

\subsection{The controlled confirmation}

This specialist supplies its own natural control, and it is the strongest
piece of evidence in the chapter. The small sibling rule
(\texttt{...|SHRED\_NAT}, $32$ clostridia rows) has its own fold. When
\emph{it} is withheld, the $620$-row clostridia rule remains in training,
so the mechanism is represented. The result:

\begin{center}\small
\begin{tabular}{lrrrr}
\toprule
Held-out rule & rows & clostridia rows left in training & $\rsq$ & bias (d) \\
\midrule
\texttt{CLOSTRIDIA|CLOSTRIDIA\_GRANA|HARD\_SEMI\_REVIEW} & 620 & 32 & $-3.860$ & $-62.2$ \\
\texttt{CLOSTRIDIA|...|SHRED\_NAT} & 32 & 620 & $\mathbf{+0.734}$ & $+1.2$ \\
\bottomrule
\end{tabular}
\end{center}

Same specialist, same model, same indicator, same products, nearly
inverted fold sizes --- and the difference in outcome tracks the
availability of the \emph{mechanism} in training, not the size of the test
set. The smaller fold is predicted well because its mechanism is
represented; the larger fold fails because its mechanism is not.

\subsection{A second, stronger control: the same rule in a different specialist}

The identical rule identifier
\texttt{V7|CLOSTRIDIA|CLOSTRIDIA\_GRANA|HARD\_SEMI\_REVIEW} also occurs in
the \emph{semi-hard} specialist, again with exactly $620$ rows. The
difference is what surrounds it: semi-hard contains three clostridia rules
($620 + 76 + 64 = 760$ rows) where hard contains two ($620 + 32 = 652$).
Withholding the same rule therefore leaves $140$ clostridia rows in
training in semi-hard, against $32$ in hard.

\begin{center}\small
\begin{tabular}{llrrrr}
\toprule
Specialist & Model & held-out rows & clostridia rows left & $\rsq$ & bias (d) \\
\midrule
\multirow{3}{*}{semi-hard/general} & Random Forest & 620 & 140 & $+0.612$ & $-3.9$ \\
 & LightGBM & 620 & 140 & $+0.672$ & $-1.4$ \\
 & XGBoost & 620 & 140 & $+0.519$ & $-10.3$ \\
\addlinespace[3pt]
\multirow{3}{*}{hard/general} & Random Forest & 620 & 32 & $-3.860$ & $-62.2$ \\
 & LightGBM & 620 & 32 & $-2.665$ & $-44.2$ \\
 & XGBoost & 620 & 32 & $-2.305$ & $-52.5$ \\
\bottomrule
\end{tabular}
\end{center}

Identical rule, identical fold size, identical models, identical protocol.
The only material difference is how much of the clostridia mechanism
survives in the training partition, and the outcome moves from
$\rsq = +0.612$ to $\rsq = -3.860$ across all three model families. It is
difficult to construct a cleaner demonstration that the governing variable
is mechanism coverage.

\begin{findingbox}[title={Verified explanation, not a hypothesis}]
Hard/general's collapse under leave-rule-out is attributable to a single
held-out rule that carries the only substantial representation of the
clostridia spoilage mechanism, whose characteristic time scale is roughly
two and a half times that of every other indicator in the specialist. The
controlled comparison above rules out fold size and product novelty as
explanations. This is the one fold in the study whose cause is established
by a within-data experiment rather than inferred.
\end{findingbox}

\fig{fig32_target_distributions}{0.98}{%
  Target distributions for the four weakest folds: the training rules
  against the held-out rule. The clostridia panel (left) shows almost
  disjoint support --- the held-out mass sits where the training data has
  almost none.}

\fig{fig33_feature_distributions_worst}{0.98}{%
  Input-feature coverage for the clostridia fold. The inputs are well
  covered by training data; it is the \emph{target} relationship that is
  not. This is the figure that distinguishes covariate shift from
  mechanism novelty.}

\tbl{tab_largest_errors}{The individual rows with the largest absolute
errors in the two worst folds, with their actual formulation context.}

\section{Fold 2: \texttt{HARD\_SEMI\_REVIEW|SHRED\_NAT}}

\begin{description}[leftmargin=1.5em, style=nextline, font=\bfseries\color{dgreen}, itemsep=1pt]
\item[Held out] $135$ rows, all \textbf{shredded cheddar}
\item[Treatments] control, chitosan, calcium propionate, lactic acid,
  protective culture, potassium sorbate
\item[Packaging] MAP at three CO\textsubscript{2} levels plus
  high-barrier sealed
\item[Result] $\rsq = -0.377$, MAE $= 7.7$ days, bias $= +6.9$ days
\end{description}

Actual targets average $38.4$ days ($s_y = 8.3$); the training partition
averages $89.7$ days. Here the held-out rule sits \emph{below} the training
range and the model over-predicts by seven days.

The distinguishing feature is \texttt{physical\_form = shredded}. Shredding
multiplies exposed surface area and shortens shelf life sharply relative to
block cheese of identical composition. Cheddar itself is well represented
in training ($179$ control rows, $58$ with chitosan, $18$ with calcium
propionate), so again the product is not novel --- the \emph{form} is
under-represented once its rule is withheld, and the model regresses
towards block-cheddar behaviour.

\section{Folds 3 and 4: the natamycin combination folds}
\label{sec:fold-natamycin}

These two folds directly test a hypothesis that was raised in advance: that
the training set may contain a cheese and a preservative separately while
containing no rows for their \emph{combination}. The data confirms it
exactly.

\subsection{\texttt{RICOTTA\_MAP|NATAMYCIN\_REVIEW} --- 8 rows}

\begin{center}\small
\begin{tabular}{lr}
\toprule
Quantity & Value \\
\midrule
Held-out rows & 8 (ricotta, fresh mass, natamycin, yeasts/moulds) \\
Ricotta rows in training & 428 \\
Natamycin rows in training & 169 \\
\textbf{Ricotta $+$ natamycin rows in training} & \textbf{0} \\
Mean $y$, training ricotta rows & 6.6 days \\
Mean $y$, training natamycin rows (18 other cheeses) & 29.0 days \\
Actual $y$, held-out rows & 9.5 days ($s_y = 1.09$, range 7.3--10.9) \\
Random Forest $\rsq$ & 0.072 \\
Random Forest MAE & \textbf{0.94 days} \\
\bottomrule
\end{tabular}
\end{center}

The model must compose two facts it has seen separately --- ricotta is
short-lived; natamycin extends shelf life --- for a combination it has
never seen. It does so almost perfectly: the average error is
\textbf{under one day}. The $\rsq$ of $0.072$ is near zero only because the
eight held-out rows span $3.6$ days in total, so $s_y^2 = 1.19$
days\textsuperscript{2} leaves essentially no variance to explain.

\subsection{\texttt{HMMC\_NAT|NATAMYCIN\_REVIEW} --- 17 rows}

The same structure with fior di latte and fresh mozzarella: $448$ and
$468$ training rows for those products, $160$ natamycin rows across
seventeen other cheeses, and \textbf{zero} rows for either product
combined with natamycin. Actual targets average $8.1$ days
($s_y = 1.83$); Random Forest achieves MAE $= 2.0$ days and
$\rsq = -0.335$.

\begin{findingbox}[title={How to present these two folds}]
These are the folds most likely to be misread, in either direction.

\emph{Do not} present $\rsq = -0.335$ as a catastrophic failure: the model
predicts an eight-day shelf life to within two days having never seen the
product--preservative combination, which is a genuinely reasonable result.

\emph{Do not} present MAE $= 0.94$ days as proof of excellent
generalization either: eight rows spanning three and a half days is far too
little evidence to support any general claim, and the near-zero $\rsq$ is
an honest signal of exactly that.

The correct statement is that on unseen product--preservative
combinations the model produces plausible compositional predictions, and
that these folds are too small to quantify how reliably.
\end{findingbox}

\section{Why the mean and the median disagree}

The soft/general specialist has median fold $\rsq = 0.875$ but mean fold
$\rsq = 0.604$. Both are correct, and the gap is entirely accounted for by
the two natamycin folds analysed above --- $8$ and $17$ rows out of
$8\,800$, together $0.28\%$ of the specialist's data, each weighted at
$1/7$ of the mean.

Hard/general is the same arithmetic with the opposite cause: median $0.630$
against mean $-0.097$, driven by the single clostridia fold, which is
\emph{not} small and where the low score is substantively meaningful.

\begin{findingbox}[title={Why this report never quotes one summary alone}]
Mean fold $\rsq$ weights an $8$-row rule equally with a $5\,654$-row rule.
Median fold $\rsq$ conceals the clostridia failure entirely. Pooled
out-of-fold $\rsq$ weights by row count and would let a small but
scientifically important rule disappear. MAE in days is robust to the
variance denominator but blind to whether the model beats a constant.

No single number is adequate. All four are reported for every specialist
throughout this document, and the recommendation in \Cref{ch:proposals} is
that they continue to be reported together.
\end{findingbox}

\section{Fold-level results for all specialists}

The three quality specialists and the three safety specialists behave
differently under this protocol, and the full per-fold tables
(\Cref{tab:tab_folds_random_forest,tab:tab_folds_lightgbm,tab:tab_folds_xgboost})
support the following summary observations:

\begin{itemize}[leftmargin=1.3em, itemsep=2pt]
\item \textbf{soft/general (7 folds).} Five folds above $\rsq = 0.87$; the
  two failures are the $8$- and $17$-row natamycin composites.
\item \textbf{semi-hard/general (12 folds).} The most stable specialist:
  every fold positive, range $0.374$--$0.955$, median $0.782$.
\item \textbf{hard/general (8 folds).} Median $0.630$, but two negative
  folds including the clostridia rule; the weakest specialist under this
  protocol.
\item \textbf{soft/safety (11 folds).} Median $0.812$, mean $0.692$; the
  spread comes from the small \emph{Listeria} composite rules.
\item \textbf{semi-hard/safety (8 folds).} Every fold positive, median
  $0.856$ --- against an original context-holdout score of $0.342$.
\item \textbf{hard/safety (10 folds).} Every fold positive, median $0.790$
  --- against an original context-holdout score of $-0.097$.
\end{itemize}

\begin{findingbox}[title={A revision to an earlier conclusion}]
The two safety specialists that looked broken under the original protocol
are not broken. Under $10$ and $8$ structured folds respectively they show
consistent, uniformly positive predictive skill (median $\rsq$ of $0.790$
and $0.856$). The original single random context split produced an
unrepresentative test partition for these specialists.

This correction should be stated explicitly at the meeting, because the
earlier reading --- that hard/safety-endpoint was fundamentally
unlearnable --- was recorded in project documentation and is wrong.
\end{findingbox}
